% !TEX root = ../main.tex
%-----------------------------------------------------------------------
%  What goes in, what comes out, and what to score the network on.
%  \input by ../main.tex -- not compiled on its own.
%-----------------------------------------------------------------------
\hd{What goes into the network.}
The network must be given enough information to identify the situation it is being asked
about. The table below is a starting point for those inputs; additions and cuts to it are
welcome.

\smallskip
\begin{tabularx}{\textwidth}{@{}Y{0.48}Y{1.52}@{}}
\toprule
\textbf{Group} & \textbf{Contents}\\
\midrule
the ocean state
  & temperature, salinity, the three velocity components and density, read from the
    forward run at the time we are asking about\\
what the state implies
  & stratification and vertical shear. Mixing acts on vertical gradients, so the
    network needs these\\
the parameters
  & the mixing coefficients, as single numbers --- vertical diffusivity is uniform in depth
    as the model is configured. Whether it belongs here at all is not yet settled: an
    ensemble spanning a range of diffusivity is
    running now to find out. If the sensitivity patterns hold their shape across that
    range, the network may generalise without it; if they shift, it has to be an input.\\
the geometry
  & depth, land--sea mask, cell size, latitude\\
what we are measuring
  & where the transport is being evaluated. The sensitivity fields I have generated so
    far are all evaluated at $26^{\circ}$N. Supplying the section as an input, rather
    than fixing it inside the network, would let us ask about sections it never saw in
    training. \textbf{That raises a question worth settling early: should the ensemble
    sample more than one section?}\\
how far back
  & the lead time, meaning how long before the measurement we are asking about\\
\bottomrule
\end{tabularx}

% There was a \pagebreak here, to keep this heading with its table. The shortened
% document no longer has room for it: forcing the break leaves the second page half
% empty and pushes the pseudocode onto a third. The heading now ends page 1 and the
% table opens page 2. Re-check the break if either section grows again.
\hd{What comes out.}
The network would return the three sensitivity fields below, one output head for each. A
single adjoint run labels all three at once, so the second and third ask nothing more of
the training set and can share the machinery that produces the first.

\smallskip
\begin{tabularx}{\textwidth}{@{}Y{0.58}Y{1.42}@{}}
\toprule
\textbf{Sensitivity to\ldots} & \textbf{What it tells us, and what shape it is}\\
\midrule
\textbf{mixing parameters}
  & how much the transport depends on vertical mixing at each point, as a 3-D map\\
\textbf{the initial ocean state}
  & how much it depends on the temperature, salinity and currents months or years
    earlier. These are 3-D maps too, so the same part of the network can produce them\\
\textbf{surface forcing}
  & how much it depends on the winds and the surface heat and freshwater fluxes, and
    when those fluxes mattered. 2-D maps, one per lead time, and probably the cheapest
    of the three to predict\\
\bottomrule
\end{tabularx}

\hd{What to score it on.}
The obvious choice is mean squared error between the predicted and the true map, but I
suspect it would fail here. On the reference run, sensitivity magnitudes cover an enormous
range --- many decades across the wet points --- and the largest values sit in the few
hundred cells nearest the section where the transport is measured. A squared error on the
raw field fits those cells and ignores the rest of the basin, so a network could score well
and still have most of the ocean wrong. My current thinking is to score it three ways
instead:

\smallskip
\begin{tabularx}{\textwidth}{@{}Y{0.52}Y{1.48}@{}}
\toprule
\textbf{Part} & \textbf{What it measures, and why it is there}\\
\midrule
\textbf{pattern} \emph{(most of the weight)}
  & whether sensitivity shows up in the right places with the right sign, scored on a
    signed-logarithmic transform of the map with its overall strength divided out, so that
    the few strongest cells cannot set the score on their own. This is the part we
    actually interpret\\
\textbf{strength}
  & whether the overall magnitude is right. Magnitude varies systematically with lead
    time; predicting it as a separate number lifts that variation out of the spatial
    problem, and the spatial problem is the hard one\\
\textbf{a direct check}
  & the two terms above treat the adjoint's output as truth, so a network could match it
    faithfully and inherit whatever the adjoint itself gets wrong. This term compares
    against the model instead: nudge a parameter in the forward run, measure how much the
    transport actually changes, and ask whether the network's prediction accounts for it.
    The perturbed runs come out of the ensemble anyway, so it is cheap\\
\bottomrule
\end{tabularx}

\smallskip
Points would be weighted by the volume of water they represent, so thin surface layers
do not count as much as the deep ocean, and land would be excluded throughout.
